\chapter{Distribution of co-expression with directional extensions} \label{ch:functions2}

\section{Overview}

A cursory reading of most descriptions of the multifunctional nature of directional extensions in Chadic languages, considering only the few examples presented to exemplify particular meanings, could give the impression that linguists have a fairly clear grasp on how these verbal extensions are used by speakers of Chadic languages. However, when broadening the scope to examine a larger selection of examples with directional extensions, a different picture emerges. Of the approximately 2,000 unique combinations of a verb and a directional extension in the dataset, in around 1,000 of those cases the meaning contributed by the verbal extension is unclear or not made explicit in the translation, as in \ref{ex:haus0}, \ref{ex:park02} and \ref{ex:molo0}. The verbal extensions in each of these examples are elsewhere shown to have directional or subsequent motion meanings. However, in these and other contexts the translation given omits any directional or subsequent motion meaning, and there is no explicit indication in the prose of what the function of the verbal extension is.

\ea\label{ex:haus0}
\langinfo{\hyperref[haus]{Hausa} verb \textit{har̃b} `shoot' with ventive suffix}{}{\cite[662]{newmanHausaLanguageEncyclopedic2000}} \\
\gll  yā har̃b-ō	\\
he shot-VENT		\\
\glt He shot (it)	

\ex \label{ex:park02}
\langinfo{\hyperref[park]{Parkwa} verb \textit{uz} `eat' with upward suffix}{}{\cite[94]{jarvisEsquisseGrammaticalePodoko1989}} \\
\gll uz-u mayə́ dafá	\\
eat.AOR-UP 1SG boule	\\
\glt je mangeai de la boule	[I ate some boule] 	

\ex \label{ex:molo0}
\langinfo{\hyperref[molo]{Moloko} verb \textit{wats} `write' with ventive extension}{}{\cite[28]{friesenGrammarMoloko2017}} \\
\gll Hʊrmbʊlɔm à-wats=ala kə ɔkʷɔr aka	\\
God 3s+pfv-write=to on stone on	\\
\glt God wrote them on the stone [tablet].
\z 

There may be several reasons for this prevalent lack of specificity regarding the function of verbal extensions in the data. It may be an issue of naturalness in translation. A verbal extension might contribute a motion-related meaning that does not translate naturally into the language of analysis (English, French or German) and so could be omitted for that reason. Alternatively, a verbal extension could be contributing a tense-aspect meaning or some other abstract grammatical function which does not have an equivalent in the target translation language, and so is not detectable unless explicitly analyzed by the linguist. A third reason why the function of an extension might be undetectable is that the linguist may incorporate the function of the extension into the gloss of the verb. Imagine, for example, that in \ref{ex:park0} the verb \textit{mbəɗ} was glossed `answer' instead of `pour'. In such a case, the function of the verbal extension would be completely obscured. 

It may also be the case that the linguist describing the grammar simply does not yet have a grasp on all of the possible functions of verbal extensions in the language. In a description of a ventive extension in \hyperref[pero]{Pero}, \citet[97]{frajzyngierGrammarPero1989} notes: ``even in simple sentences the functions as described above are not always obvious, or maybe it is the Indo-European background of this linguist that prevents the proper understanding of the choice of ventive rather than non-ventive variants...''

Whatever the reasons for so many examples of verbal extensions with an unclear function, it should be kept in mind that there are many examples that cannot be included in the analysis due to a lack of clarity about the function of the directional extension. Among the examples of directional extensions found in non-motion contexts, only about half of the available data specifies what the function of the extension is. It is unknown whether the other half would exhibit similar patterns if more precise translations were made available, or if the lack of specificity in the unclear cases is due to those extensions having more subtle functions that have not yet been consistently described across Chadic languages. 

Table \ref{tab:DIR_plus} provides an overview of how frequently various meanings are co-expressed with types of directional meaning across Chadic languages. The first column (DIR) repeats from Table \ref{tab:directionals} in  Chapter \ref{ch:distribution} regarding how many of the 91 languages in the data have a directional extension of each type of path (ventive, itive, etc.). The second column (+AM/LOC) indicates how many languages show a pattern of co-expressing an associated motion or locational meaning in the same verbal extension that expresses a directional meaning. This means that in the 46 languages with a ventive directional extension, that morpheme can also express a motion or location meaning. In contract, in the 19 languages with an itive directional extension, only seven can also express a motion or location meaning. 

The third column (+Aspect) shows how frequently aspectual and other adverbial meanings are co-expressed with each directional meaning. The fourth column (+Idiom) counts languages where the directional extension is attested in idiomatic uses. The fifth column (+CAM) tabulates co-expression of caused accompanied motion. The final column (+Unclear) gives the number of languages in which a directional extension is also attested in examples where the meaning contributed by the extension is not discernible in the available data. Note in the final column that the issue of an indiscernible function of a verbal extension is widespread across languages and paths of motion.\footnote{In some cases, the lack of any ``unclear'' examples is due to a limited amount of data, rather than an abundance of exceptionally explicit translations.} 

%data in tab:DR_plustable2 is from Python code in Functions 
\begin{table}
	\caption{Number of languages with other meanings co-expressed in a directional extension}
	\label{tab:DIR_plus}
	\begin{tabular}{lllllll} % add l for every additional column or remove as necessary
		\lsptoprule
		{} &  DIR &  +AM/LOC &  +Aspect &  +Idiom &  +CAM &  +Unclear \\
		\midrule
		VENT &   46 &                40 &                 6 &             7 &    25 &        29 \\
		ITIVE &   19 &                 7 &                 0 &             9 &    11 &        14 \\
		INTO &   16 &                 4 &                 2 &             3 &     5 &        11 \\
		OUT &   15 &                 3 &                 4 &             7 &    11 &        15 \\
		UP &    8 &                 5 &                 1 &             5 &     5 &         8 \\
		DOWN &   13 &                 4 &                 1 &             3 &    10 &        11 \\
		ONTO &    9 &                 4 &                 4 &             6 &     0 &         7 \\
		UNDER &    3 &                 2 &                 0 &             2 &     0 &         3 \\
		\lspbottomrule
	\end{tabular}
\end{table}


\section{Motion and location} \label{sec:comotion2}

%data in tab:DR_plustable is from Python code in Functions 
%\begin{table}
%	\caption{Directional extensions with additional functions co-expressed}
%	\label{tab:DR_plusTABLE}
%	\begin{tabular}{llllll} % add l for every additional column or remove as necessary
	%		\lsptoprule
	%		{} &  DIR &  +PriorAM &  +ConcAM &  +Subs(C)AM &  +LOC \\
	%		\midrule
	%VENT &   49 &        4 &       1 &         29 &    5 \\
	%ITIVE &   22 &        0 &       0 &          6 &    1 \\
	%INTO &   17 &        0 &       0 &          1 &    1 \\
	%OUT &   17 &        0 &       0 &          2 &    0 \\
	%UP &   12 &        0 &       1 &          3 &    0 \\
	%DOWN &   16 &        0 &       0 &          3 &    0 \\
	%ONTO &   11 &        0 &       0 &          1 &    1 \\
	%UNDER &    3 &        0 &       0 &          0 &    2 \\
	%VENT.OUT &    1 &        0 &       0 &          1 &    0 \\
	%ITIVE.DOWN &    1 &        0 &       0 &          0 &    0 \\
	%		\lspbottomrule
	%	\end{tabular}
%\end{table}

Table \ref{tab:DR_plusTABLE2} shows the number of languages that have a verbal extension in which directional meaning is co-expressed with a motion or location meaning. The column DIR shows the number of languages that have a directional extension of each of the most common types. The other columns indicate in how many languages a directional extension of that type co-expresses another meaning: Prior associated motion (+PriorAM), concurrent associated motion (+ConcAM), subsequent associated motion or subsequent caused accompanied motion (+Subs(C)AM), and location (+LOC).

%data in tab:DR_plustable2 is from Python code in Functions 
\begin{table}
	\caption{Directional extensions and motion/location co-expression}
	\label{tab:DR_plusTABLE2}
	\begin{tabular}{llllll} % add l for every additional column or remove as necessary
		\lsptoprule
		{} &  DIR &  +PriorAM &  +ConcAM &  +Subs(C)AM &  +LOC \\
		\midrule
		VENT &   46 &        3 &       0 &         30 &    7 \\
		ITIVE &   19 &        0 &       0 &          7 &    0 \\
		INTO &   16 &        0 &       0 &          1 &    3 \\
		OUT &   15 &        0 &       0 &          3 &    0 \\
		UP &    8 &        0 &       1 &          4 &    0 \\
		DOWN &   13 &        0 &       0 &          4 &    0 \\
		ONTO &    9 &        0 &       0 &          1 &    3 \\
		UNDER &    3 &        0 &       0 &          0 &    2 \\
		\lspbottomrule
	\end{tabular}
\end{table}

As is clear from Table \ref{tab:DIR_plus}, verbal extensions that express ventive directional meaning are the most likely to co-express another motion or location meaning. In 32 of the 46 languages (70\%) that have a verbal extension that expresses ventive directional meaning, the same form is also reported to have another motion or location meaning.\footnote{The \hyperref[baca]{Bacama} and \hyperref[mbud]{Mbudum} ventive extensions co-express both prior and subsequent motion. The \hyperref[muya]{Muyang}, \hyperref[haus]{Hausa}, \hyperref[bole]{Bole} and \hyperref[park]{Parkwa} ventive extensions co-express both subsequent motion and locative. The \hyperref[kana]{Kanakuru} ventive forms co-express prior motion, subsequent motion and locative.} Table \ref{tab:DR_plusTABLE2} reveals that the most commonly found meaning co-expressed with directional meaning is subsequent motion. It is comparatively rare to find other types of motion or location meaning co-expressed with directional meaning.

\subsection{Co-expression with ventive directional extensions} \label{sec:covent2}

As summarized above in Table \ref{tab:DR_plusTABLE2}, the majority of ventive directional extensions are described as being co-expressive with other motion and location meanings. Table \ref{tab:vent_plusTABLE} shows the distribution of types of co-expressed meaning across the four branches of Chadic languages. Co-expression of subsequent motion with ventive direction is found to be common across all branches except for the three Masa languages. The rare cases of co-expression with prior associated motion and locative meaning are attested in Biu-Mandara and West Chadic languages.

%Data in table vent_plusTABLE is tabulated in Python file FunctionsBranch
\begin{table}
	\caption{Ventive directional extensions and co-expression by branch}
	\label{tab:vent_plusTABLE}
	\begin{tabular}{llllll} % add l for every additional column or remove as necessary
		\lsptoprule
		{} &  Ventive &  +PriorAM &  +ConcAM &  +Subs(C)AM &  +LOC \\
		\midrule
		West &       14 &         1 &        0 &          10 &     4 \\
		Biu-Mandara &       26 &         2 &        0 &          18 &     3 \\
		Masa &        3 &         0 &        0 &           0 &     0 \\
		East &        3 &         0 &        0 &           2 &     0 \\
		Total &       46 &         3 &        0 &          30 &     7 \\
		\lspbottomrule
	\end{tabular}
\end{table}

\subsection{Co-expression with itive directional extensions} \label{sec:coitive2}

As discussed in Chapter \ref{ch:distribution}, itive directional extensions are less common than ventive directional extensions. They are primarily limited to Central Chadic languages with one exception. Table \ref{tab:itive_plusTABLE} shows that of the relevant types of meaning only subsequent motion is found to be co-expressed with itive directional meaning in the same verbal extension, and only in Biu-Mandara languages.\footnote{There is one marginal case in \hyperref[mada]{Mada} of a locative interpretation of a form also used for itive direction. The verb \textit{grá} `do' in a palatalized form has an altrilocative interpretation which is analogous to the (palatalized) itive form of the verb \textit{ɬé} `go home' \citep[29]{ernstkurdiTenseAspectMood2016}.}

%Data in table vent_plusTABLE is tabulated in Python file FunctionsBranch
\begin{table}
	\caption{Itive directional extensions and co-expression by branch}
	\label{tab:itive_plusTABLE}
	\begin{tabular}{llllll} % add l for every additional column or remove as necessary
		\lsptoprule
		{} &  Itive &  +PriorAM &  +ConcAM &  +Subs(C)AM &  +LOC \\
		\midrule
		West &      0 &         0 &        0 &           0 &     0 \\
		Biu-Mandara &     16 &         0 &        0 &           7 &     0 \\
		Masa &      2 &         0 &        0 &           0 &     0 \\
		East &      1 &         0 &        0 &           0 &     0 \\
		\midrule
		Total &     19 &         0 &        0 &           7 &     0 \\
		\lspbottomrule
	\end{tabular}
\end{table}



\section{Aspectual and other adverbial meanings} \label{sec:cogrammatical2}
%numbers in this section are from the Python file Functions 

Table \ref{tab:DIR_gram} gives a quantitative overview of how many languages show evidence of one of these co-expressed functions. These include a completive or totality function (COMPL) and other tense and aspect (TAM, Section \ref{sec:tense}), marking a benefactive (BEN, Section \ref{sec:beneficiary}), additive meaning (ADD, Section \ref{sec:additive}) and marking the object as affected by the action such that it splits into parts (SPLIT, Section \ref{sec:objectsplit}).

%data in tab:DIR_gram from python script in Functions
\begin{table}
	\caption{Directional extensions and aspectual/adverbial uses}
	\label{tab:DIR_gram}
	\begin{tabular}{lllllll} % add l for every additional column or remove as necessary
		\lsptoprule
		Path &  DIR &  COMPL &  TAM &  BEN &  ADD &  SPLIT \\
		\midrule
		VENT &   46 &      0 &    1 &    5 &    0 &      0 \\
		ITIVE &   19 &      0 &    0 &    0 &    0 &      0 \\
		INTO &   16 &      0 &    0 &    0 &    0 &      2 \\
		OUT &   15 &      3 &    0 &    1 &    0 &      0 \\
		UP &    8 &      1 &    0 &    0 &    0 &      0 \\
		DOWN &   13 &      1 &    0 &    0 &    0 &      0 \\
		ONTO &    9 &      0 &    0 &    0 &    4 &      0 \\
		BUSH &    1 &      1 &    0 &    0 &    0 &      0 \\
		\lspbottomrule
	\end{tabular}
\end{table}